\scanpage{113}

\deleted{Donc par définition de l'iso du présent \S\ $1'$, appliqué à $(\mathcal M,q_{\mathcal M},i\omega_{\mathcal M})$, on l'obtient par transport de structure}

On doit avoir évidemment

\[
f_{q,\omega}=\theta_{q,\omega}\,g_q
\]

où $\theta_{q,\omega}$ est un autom de $\Sigma$, qui est compatible à tt chgt de base, et fonctoriel en $(E,q,\omega)$. Donc pour prouver $\theta_{q,\omega}=\operatorname{id}$ OPS que $S=\operatorname{Spec}\mathbb Z[\frac12]$, $(E,q,\omega)$ étant la forme \unsure{standart|standard}. On a une section $a_0$ de $\Sigma$, et $\Sigma$ s'identifie au quotient $G/T$, où $G={\unsure{\operatorname{GP}(1)|\operatorname{GP}(A)|\operatorname{SO}(3)}}_S=\operatorname{Aut}(E,q,\omega)$, et $N$ est le stabilisateur de $a_0$ i.e. un tore maximal. Le schém \deleted{\unsure{seul|seules}}\inserted{en groupes des} autom de $\Sigma\simeq G/T$ commut$^{\mathrm t}$ à $G$ est $N(T)/T\simeq(\mathbb Z/2\mathbb Z)_S$, $\mathbb Z/2\mathbb Z$ opérant précisément par l'antipodisme.

Ceci montre que $\theta_{q,\omega}$, qui est donné par une section de $N(T)/T$ sur $S$, est soit l'identité, soit l'antipodisme. Mais lequel des deux ?

% archive sheet number: 112
